% ================================================================
% SECTION 4 — BAILOUT PROTOCOL: IMPLEMENTATION IN THE BX3 FRAMEWORK
% ================================================================
\section{Implementation in the BX3 Framework}\label{sec:bailout-implementation}

The full formal specification of the Bailout Protocol is given in \S~\ref{sec:bailout-full}.
This section describes how that specification is embedded in the BX3 three-layer
architecture, detailing the implementation of each protocol component across the
Purpose, Bounds, and Fact layers, and establishing the precise relationship
between the Bailout Protocol and the upstream accountability guarantee that
BX3 systems fail upward into human consciousness, never downward into autonomous chaos.

% ---------------------------------------------------------------
\subsection{Purpose Layer: The Human Accountability Anchor}\label{sec:purpose-anchor}

The Purpose layer is the structural terminus for every Bailout Protocol escalation.
It is not a passive recipient of exceptions — it is the definitive accountability
anchor, carrying the human principal's intent and serving as the only permissible
termination point for unresolved exception states.

At node provisioning, the Purpose layer records a reference to a named human
principal $h(n)$ and stores it as an immutable, Fact-layer-protected attribute.
This is not a configuration parameter — it is a structural commitment.  The
human anchor is not selectable at runtime; it is fixed at provisioning and
cannot be changed by any machine actor.  This immutability is enforced by the
Fact layer's pre-commit hook (IC-4): any proposed modification to $h(n)$ requires
a new ledger entry authorized by the current $h(n)$ itself, creating a circular
authorization constraint that renders the field effectively immutable in
practice and structurally immutable by architecture.

The Purpose layer carries two responsibilities relevant to the Bailout Protocol.
First, it defines the operating range $R(n)$ for each node — the provisioned
boundary within which the Bounds Engine may operate autonomously.  This
definition is the reference against which the bail-out trigger condition TC
(\S~\ref{sec:bailout-trigger}) is evaluated: an exception $x$ is flagged as
potentially bail-out-eligible precisely when $x \notin R(n)$.  The Purpose layer
does not participate in the trigger evaluation — it defines the boundary that
makes the evaluation possible.

Second, the Purpose layer receives escalated exceptions via the escalation
algorithm (\S~\ref{sec:escalation-algorithm}) and issues binding resolution
directives.  These directives are recorded in the authorization ledger via the
Fact layer before they take effect, satisfying Invariant~\ref{inv:intent} and
Invariant~\ref{inv:ledger-immutability}.  The only permissible directives are
\texttt{ACK} (acknowledge and continue), \texttt{RESOLVE} (execute a specific
binding action), and \texttt{REJECT} (disallow the proposed action and enter
a quiescent error state).  No machine actor may issue any of these directives;
the protocol requires a human principal to supply the binding intent.

The Purpose layer enforces what may be called the \emph{non-substitutability
property}: no machine actor operating under delegated authority possesses the
principal's intent in sufficient completeness to substitute for $h(n)$ in the
accountability chain.  This is not a limitation of current AI systems — it is a
structural argument.  The principal's intent is defined at provisioning and
carried exclusively in the Purpose layer; any machine actor's authority is
inherently a strict subset of that intent, operating within $R(n)$, and
therefore cannot serve as a terminus for exceptions that exceed its own scope.
The Bailout Protocol makes this structural argument operationally enforceable:
exceptions that exceed $R(n)$ have no machine-actor terminus available to them
except through deliberate protocol violation.

% ---------------------------------------------------------------
\subsection{Bounds Engine: Triggering Bail-Out}\label{sec:bounds-trigger}

The Bounds Engine is the operative component of the Bounds layer and the
direct triggering authority for the Bailout Protocol.  It performs two functions
in the bail-out mechanism: continuous monitoring against $R(n)$, and deterministic
signal emission when the trigger condition is satisfied.

The Bounds Engine maintains a read-only view of $R(n)$, provisioned at node
startup and protected against runtime modification by the Fact layer's
pre-commit hook.  On every consequential action proposed by the node, the
Bounds Engine evaluates whether the action falls within $R(n)$ by querying the
authorization ledger via the Fact layer (IC-1).  If the action is within
$R(n)$, it proceeds through the standard execution path.  If the action — or
an observed exception condition $x$ — falls outside $R(n)$, the Bounds Engine
evaluates the full trigger condition TC:

\[
\text{TRIGGER}(n, x) \equiv x \notin R(n) \;\land\; \neg\text{resolved}(n, x) \;\land\; \text{confidence}(n, x) < \tau .
\]

All three sub-conditions must be satisfied simultaneously.  The first
sub-condition ($x \notin R(n)$) is a deterministic query against the provisioned
operating range.  The second ($\neg\text{resolved}(n, x)$) is a lookup against
the authorization ledger to confirm that $x$ has not already received human
adjudication in the current session.  The third ($\text{confidence}(n, x) < \tau$)
is a provisioned threshold — immutable at runtime — that prevents the Bounds
Engine from inflating its own confidence to suppress warranted escalation.

If TC evaluates to true, the Bounds Engine transitions the protocol state
machine from \text{IDLE} to \text{BOUNDED} (T1) and begins a time-bounded
resolution attempt.  The resolution window is governed by $T_{\text{bounds}}$,
enforced by the Fact layer's pre-commit hook rather than the Bounds Engine.
This architectural placement is deliberate: if $T_{\text{bounds}}$ were
enforced by the Bounds Engine, a compromised or misconfigured machine actor
could suppress the timeout to prevent transition to \text{BAIL\_PENDING}.
By enforcing the timeout at the Fact layer, the framework ensures that
timeout expiration is structurally compelled, not operationally optional.

The Bounds Engine also emits an explicit \texttt{bailout\_signal} independently
of TC when it detects a logical contradiction within $R(n)$ — a condition
where two individually valid inputs cannot be jointly resolved.  This secondary
trigger bypasses TC entirely and initiates the transition to \text{BAIL\_PENDING}
via T3, ensuring that structural irresolvability within $R(n)$ is treated as
equivalent to a scope-exceeding exception for purposes of the bail-out mechanism.

The no-discretion property (TC-secondary) ensures that the trigger evaluation
is architecturally enforced: there is no machine-actor gate between trigger
evaluation and protocol activation.  The Bounds Engine does not choose whether
to evaluate TC — it is a read-only query against provisioned data structures,
and the result directly drives the state machine.  This eliminates the possibility
that a machine actor could suppress a warranted escalation by deferring or
misrepresenting the trigger evaluation.

% ---------------------------------------------------------------
\subsection{Fact Layer: Enforcing the Audit Trail and State Machine}\label{sec:fact-enforcement}

The Fact layer is the architectural foundation of the Bailout Protocol's
accountability properties.  It enforces three distinct functions: the
authorization ledger that gates all consequential actions, the Forensic Ledger
that records every protocol event, and the state machine transition relation
that governs the protocol's runtime behavior.

The state machine $\mathcal{B} = (Q, q_0, \Sigma, \rightarrow, Q_F)$ is embedded
in the Fact layer of every BX3 node.  This is a design choice with architectural
consequences: by placing the state machine in the Fact layer, the framework
ensures that the protocol's transitions are enforced by the same mechanism that
enforces all other BX3 invariants — the Fact layer's pre-commit hook.  No Bounds
layer component, no parent node, and no external agent may modify, stall, or
bypass a state transition once its preconditions are satisfied.  The state
machine is therefore not merely a formal model — it is a structurally operative
constraint on every node's execution.

The deterministic transition relation (Table~\ref{tab:bailout-transition}) is
enforced by the Bailout Monitor, a Fact-layer component that evaluates all
transition preconditions on every event and commits at most one transition per
event to the Forensic Ledger.  The priority function $\pi(n, x)$ (Equation PF)
resolves simultaneous trigger conditions in a deterministic total order,
ensuring that the protocol's behavior is reproducible and auditable.  The
priority function's weights $w_1, w_2, w_3$ are provisioned at node initialization
and are equally immutable at runtime — no machine actor may adjust the ordering
of trigger conditions to suppress high-priority escalations.

The bypass mechanism (Definition~\ref{def:bypass}) is also Fact-layer enforced.
The Fact layer's pre-commit hook rejects any ledger entry for a resolution
where the trigger condition is satisfied and $T_{\text{bounds}}$ has expired.
This prevents the Bounds Engine from issuing a retroactive resolution to
suppress a warranted escalation after the timeout has fired.  The parent pointer
$\alpha(n)$ and human-root identifier $h(n)$ are stored in Fact-layer-managed
memory with read-only access at the machine-actor level.  A machine actor cannot
redirect the parent chain, insert itself as an intermediate resolver, or modify
the human-root field that terminates the chain.  Any proposed modification
requires Purpose layer authorization via a new ledger entry — which cannot be
completed without human approval if the modification would affect the bailout
terminus.  The bypass is therefore not merely a routing constraint — it is a
structural property enforced at the pre-commit stage by the Fact layer.

The Forensic Ledger records every state transition, every trigger condition
evaluation, every timeout event, every pathway selection decision, and every
human directive issued by $h(n)$.  Each entry is a 6-tuple:

\[
f = (\text{id}(f),\; t_{\text{logical}}(f),\; \text{node}(f),\; \text{event}(f),\; \text{context}(f),\; \text{hash}(f))
\]

where $\text{hash}(f)$ is the SHA-256 hash of the preceding entry, forming a
cryptographically chained sequence from the genesis entry to the most recent.
The cryptographic chaining ensures that no entry can be modified, deleted, or
inserted without detection — a property verifiable by any auditor without
access to the ledger's contents.

The Forensic Ledger serves two functions in the Bailout Protocol's implementation.
First, it is the authoritative source of truth for the runtime state of the
protocol state machine at every node.  Before any state transition is committed,
the Fact layer queries the Forensic Ledger to confirm the current state and to
record the transition.  Second, it is the permanent evidence record that
demonstrates to any external auditor — human principal, regulator, or third-party
reviewer — that the protocol operated correctly: that every exception reached
$h(n)$, that no machine actor absorbed or suppressed any escalation, and that
all diagnostic context was preserved and attributed at every step.

% ---------------------------------------------------------------
\subsection{Relationship to the Upstream Accountability Guarantee}\label{sec:bailout-accountability}

The upstream accountability guarantee — that BX3 systems fail upward into
human consciousness, never downward into autonomous chaos — is realized through
the Bailout Protocol as the runtime expression of that guarantee.  The guarantee
is not a policy aspiration; it is a structural property enforced by three
independent mechanisms: the Purpose layer's non-substitutability property,
the Bounds Engine's no-discretion trigger, and the Fact layer's pre-commit
enforcement of the state machine and bypass mechanism.

To see why the Bailout Protocol is the runtime expression of the guarantee,
consider the failure mode it prevents.  In any multi-agent system without an
architectural bail-out mechanism, an exception encountered by a node may be
handled at any level in the agent hierarchy: resolved locally, suppressed by a
parent node, absorbed by an intermediate tier, or simply not propagated because
no propagation pathway exists.  Each of these outcomes violates the upstream
accountability guarantee by disconnecting the exception from the human principal
who bears ultimate responsibility for the system's actions.

The Bailout Protocol eliminates these failure modes structurally.  The trigger
condition TC ensures that any exception outside $R(n)$ is automatically flagged,
without machine-actor discretion.  The state machine ensures that the flag is
not suppressible: once in \text{BOUNDED}, the node must either produce a
Fact-layer-approved resolution within $T_{\text{bounds}}$ or transition to
\text{BAIL\_PENDING} and then to \text{ESCALATING}, with the latter transition
occurring atomically via the Fact layer without an intervening machine-actor
decision step (T4).  The bypass mechanism ensures that the escalation reaches
$h(n)$ and no other actor.  The Forensic Ledger ensures that the entire chain
is permanently attributable.

Together with the other four pillars — Loop Isolation (preventing state
circularity that would make exception provenance untraceable), Recursive Spawning
(maintaining the parent-child accountability chain that provides the propagation
substrate), Spatial Firewall (preventing cross-domain contamination that would
obscure exception attribution), and Sandbox Gate (enforcing the pre-commit hooks
that make Bounds Engine operating range enforcement structurally operative) —
the Bailout Protocol makes the upstream accountability guarantee architecturally
compelling rather than operationally aspirational.

The relationship between the five pillars and the upstream accountability
guarantee is one of mutual reinforcement.  Loop Isolation ensures that the
provenance of every exception is traceable; Recursive Spawning ensures that
the propagation path to $h(n)$ is structurally available; Spatial Firewall
ensures that exceptions are cleanly attributed to their originating domain;
Sandbox Gate ensures that inter-domain communication does not bypass the
pre-commit hooks that enforce the Bailout Protocol's trigger conditions; and
the Bailout Protocol ensures that every unresolved exception reaches $h(n)$
with forensic attribution.  No pillar is sufficient alone; all five are required
to make the guarantee unconditional.

% ---------------------------------------------------------------
\subsection{Forensic Ledger Recording}\label{sec:forensic-recording}

The Forensic Ledger's role in the Bailout Protocol extends beyond passive
recording to active enforcement of protocol correctness.  Three categories of
events are recorded with sufficient detail to support both runtime state
verification and retrospective audit.

\textbf{State transition records.} Every transition of the protocol state machine
$\mathcal{B}$ is recorded in the Forensic Ledger before any subsequent action
is taken by any layer of the node.  The record includes the originating state,
the triggering event, the target state, the logical timestamp assigned by the
Fact layer, and the node identifier.  This record satisfies Invariant~INV$_1$:
the complete state history of the protocol at each node is reconstructable from
the ledger, enabling verification that no transition was suppressed, reordered,
or incorrectly recorded.

\textbf{Escalation propagation records.} Each invocation of the \texttt{SendToParent}
function during the escalation algorithm records: the sending node identifier,
the exception record (type, severity, origin timestamp), the diagnostic context
chain accumulated at each intermediate node, the selected communication pathway,
the timeout value at time of send, and the acknowledgment result.  These
records enable reconstruction of the complete escalation path and confirm that
no intermediate node intercepted, redirected, or absorbed the exception — a
necessary condition for verifying the bypass property (BP).

\textbf{Human directive records.} Every directive issued by $h(n)$ — \texttt{ACK},
\texttt{RESOLVE}, or \texttt{REJECT} — is recorded with the full directive content,
the logical timestamp, the originating node, and theFact layer hash of the
triggering state.  The hash provides cryptographic integrity verification: any
subsequent modification to the triggering state that would change the basis for
the directive would break the chain and be detectable by any auditor.  This
record satisfies Invariant~INV$_3$: every RESOLVED state is terminal with a
recorded directive from $h(n)$.

The Forensic Ledger is append-only: no machine actor may modify, delete, or
suppress any entry.  This immutability is structurally enforced by the Fact
layer's pre-commit hook, which rejects any proposed write operation that would
alter a committed entry.  The SHA-256 chaining ensures that tampering with any
entry — whether by modification, deletion, or insertion — breaks the chain at
the point of tampering, making the break detectable without access to the
ledger's contents.  An auditor need only verify the chain integrity by
recomputing each hash and comparing it to the stored successor hash to confirm
that no tampering has occurred.

% ---------------------------------------------------------------
\subsection{State Machine as Fact Layer Extension}\label{sec:state-machine-extension}

The protocol state machine $\mathcal{B}$ is formally described as a component
of the Fact layer — not as a separate or auxiliary mechanism.  This placement
has four architectural consequences that directly support the upstream
accountability guarantee.

\textbf{Uniform enforcement.} Because the state machine is embedded in the
Fact layer, its transition relation is enforced by the same pre-commit hook
that enforces all other BX3 invariants (IC-1 through IC-4).  There is no
separate enforcement path for the Bailout Protocol — the Fact layer pre-commit
hook evaluates all transition preconditions and commits all state changes
according to a single, consistent enforcement logic.  This eliminates the risk
of divergent enforcement paths that could allow the protocol's constraints to
be bypassed through a different mechanism than the one that enforces the
framework's other invariants.

\textbf{Precedence ordering.} The Bailout Monitor's priority function $\pi(n, x)$
establishes a deterministic total ordering over simultaneous trigger conditions.
When multiple exceptions fire concurrently at a node — a condition that can
occur in heavily parallelized agent graphs — the priority function selects the
highest-priority trigger and commits its transition before any other.
The priority ordering is provisioned at node initialization and is immutable
at runtime, preventing a machine actor from manipulating the ordering to suppress
high-severity escalations in favor of lower-severity ones that are easier to
resolve autonomously.

\textbf{Cryptographic state binding.} The Fact layer hash of the triggering state,
included in every human directive record, cryptographically binds the directive
to the exception context that prompted it.  This prevents a class of attacks in
which a machine actor modifies the node's state after an exception is observed
but before the human principal receives the escalation, thereby presenting a
misleading state to $h(n)$ that would cause an incorrect directive to be issued.
The hash chain in the Forensic Ledger makes any such modification detectable,
and the directive record's hash field makes the specific tampering attributable
to the node and time window in which it occurred.

\textbf{State isolation.} The atomic state update property — no intermediate or
partially-executing states are observable externally — ensures that the protocol's
internal logic cannot be inspected or manipulated by a parent node or external
observer during a transition.  A parent node observes only the committed state
values in the Fact layer worksheet; it cannot observe or interfere with the
transition logic as it executes.  This prevents a parent node from stall-hacking
the protocol by observing an in-flight transition and issuing a competing
transition that could corrupt the state machine's deterministic execution.

The state machine extension to the Fact layer also provides the implementation
substrate for the protocol's liveness guarantees.  The timeout parameters
$T_{\text{bounds}}$, $T_{\text{prop}}$, $T_{\text{cap}}$, and $T_{\text{human}}$
are enforced by the Fact layer's scheduling mechanism, not by the Bounds Engine
or any machine-actor component.  This placement ensures that timeout expiration
is structurally compelled and cannot be suppressed, artificially extended, or
misrepresented by any machine actor.  The Pathway Health Registry, which tracks
the availability and latency of communication pathways to $h(n)$, is also a
Fact-layer-managed runtime data structure updated by a background health-ping
thread operating at the same scheduling priority as the Bailout Monitor.  Its
state is observable via the Forensic Ledger and is therefore auditable.

The combined effect of these Fact-layer enforcement mechanisms is that the
Bailout Protocol's guarantees are not dependent on the correct behavior of
any machine actor.  The protocol operates correctly regardless of whether the
Bounds Engine is compromised, whether a parent node is unresponsive, or whether
an intermediate machine actor attempts to intercept the escalation.  The Fact
layer is the fail-safe anchor: it is the only layer that may issue a definitive
halt to autonomous action, and it does so by refusing to approve any proposed
action that lacks a valid authorization entry.  The Bailout Protocol is the
specific operationalization of this fail-safe function for unresolved exception
states — and the Fact layer is the structural guarantee that the operationalization
is unconditional.